\documentclass{article}
\usepackage{amsmath}
\begin{document}
\title{An unrelated render wake can supply or consume a qualified capture}
\maketitle

\paragraph{Scope.}
The source semantics are those of the pinned Gazebo3 sensor-system caller
and the pinned Sensors3 manager and rate gate. The caller's CMake requires
Sensors3 at version 3.3 or later; this compatibility evidence does not
identify the binary in a deployed image. The replay executes the selected
methods in a serialized, initialized fixture with fixed, positive-rate
members, at distinct increasing simulation times. Scene update is stubbed;
a renderable instrument's callback records the supplied time. There is no
wall-clock, event transport, concurrent interleaving, or real image claim.
The service predicate consumes an exposure taken at or after the beginning
of its qualified sensing window, not the latest cached image as such.

\paragraph{Two distinct gates.}
\texttt{PostUpdate} uses a mask and a nominal next update date to decide
which sensor can wake the render worker. Mask expiry at equality removes
the mask. When at least one sensor is active, \texttt{RunOnce} renews the
masks of only those active sensors and calls the ordinary manager for all
members. The manager uses the nominal due date without that wake mask.
A pending-scene-only wake can update the scene without running this
manager call. A regular update advances the previous due date by one
period, even if that date was already in the past. A sensor's own mask
is therefore neither a guarantee of no capture nor a guarantee that its
next eligible capture will be serviced.

\paragraph{Exact count on a masked interval.}
Fix a time interval strictly inside target \(A\)'s mask, initial nominal
due date \(d\), period \(P>0\), and a sorted set \(W\) of distinct
external manager-wake times on the interval. A counted wake is a pass
with a nonempty active set, not a scene-only wake. The assumptions
exclude forces, resets, membership changes, and a change of \(W\) caused
by a target capture. These times must be known or conservatively related
by source evidence for operational use; they are not a controller
oracle. For a prefix ending at \(t\), write \(W_t=W\cap(-\infty,t]\).
Then the number of target captures is exactly
\[
 N(t)=\min_{0\le k\le |W_t|}
   \left(k+\#\{w\in W_t:w\ge d+kP\}\right).
\]
At most \(k\) captures can precede \(d+kP\); each later capture
uses a distinct wake. To see tightness, run the rate gate greedily in
time order. If no wake is skipped, the \(k=0\) cut is tight.
Otherwise let \(k\) be the number of captures before the last skipped
wake. That wake is earlier than \(d+kP\), as are all its predecessors.
All successors are used and occur at or after \(d+kP\). This cut is
tight. For a qualified window \([q,t]\), at least one target update
is possible in the replayed schedule exactly when \(N(t)>N(q^-)\).
The cumulative count is monotone under adding wakes, but the
difference for a specified acquisition window need not be. This is
nested interval matching on an explicitly projected wake schedule,
not a speedup over an equally informed scheduler.

\paragraph{Opposite effects of extra activity.}
In milliseconds relative to a setup tick, \(A\) starts nominally due
at \(-48\), fires at \(0\), has its mask through time \(45\)
(exclusive), and is now due at \(2\). In the tested relative window,
old sensor \(B\) becomes due at \(11\) or is not yet due. Without an
outside wake, a pending scene update cannot capture \(A\) during
the mask. When \(B\) fires at \(11\), it yields its own old frame and
also an \(A\) frame after the qualified opening \(q=10\).
If an optional sensor \(C\) also becomes due at \(9\), that earlier
wake consumes \(A\)'s scheduled slot and advances its nominal due
date to \(52\). \(B\) still fires at \(11\); no usable \(A\) frame
is captured then. Thus the same old frame survives the change and
does not certify the target capture. Moving an unrelated configuration
change earlier can hurt even when the target configuration and
qualified deadline are unchanged.

In the prior far-site stop-and-go motion model, the capture at
\(q+1\) ms is sufficient with the declared 105 ms decode delay:
\(\epsilon+10(L-0.1)^2=0.00125 < 7-6.37-4/6.37\).
An exposure at \(q-1\) ms is not eligible; after the early wake
the next possible regular exposure is no earlier than \(q+42\) ms,
even with further permitted wakes. That misses the same motion
bound even with zero processing delay. This joins source
simulation timestamps to the model by a translation of the clock
origin. It assumes the declared capture/processing contract and does
not measure or enforce it. The direct scheduling baseline with the
same external wake sequence and deadlines agrees.

\paragraph{Retained evidence and residuals.}
An unrelated wake reservation is a usable recovery resource if its
time and scope are qualified and earlier wakes that can consume the
same due slot are accounted for. Retaining only an active list, a
target rate, or the old sensor's successful frame would omit a
material distinction. The scope also excludes scene/clock rewinds,
render completion latency, and actual packet delivery. No
independent mission was used.
\end{document}
